% Figure: absorption edge of a multiple-quantum-well electro-absorption
% modulator under two reverse-bias fields. A field applied across the wells
% shifts the excitonic absorption edge to longer wavelength (the quantum-confined
% Stark effect), so at a fixed operating wavelength the material switches from
% transparent to absorbing as the bias is applied. All points entered explicitly.
\begin{figure}[htbp]
\centering
\begin{tikzpicture}[font=\scriptsize]
\begin{axis}[
    width=10.5cm, height=6.2cm,
    xlabel={wavelength $\lambda$ (\si{\nano\meter})},
    ylabel={absorption coefficient $\alpha$ (\si{\per\micro\meter})},
    xmin=1490, xmax=1580,
    ymin=0, ymax=1.2,
    xtick={1490,1510,1530,1550,1570},
    ytick={0,0.3,0.6,0.9,1.2},
    grid=both,
    grid style={enggray!25},
    tick label style={font=\tiny},
    label style={font=\scriptsize},
    axis line style={enggray},
    legend style={font=\tiny, at={(0.03,0.97)}, anchor=north west,
      draw=enggray, fill=white},
    every axis plot/.append style={very thick},
]
% zero-field absorption edge (sharp exciton near 1520 nm)
\addplot[engblue] coordinates {
  (1490,1.05) (1500,1.00) (1510,0.90) (1518,0.75) (1524,0.45)
  (1528,0.20) (1532,0.07) (1538,0.02) (1545,0.008) (1555,0.004)
  (1570,0.002) (1580,0.001)
};
\addlegendentry{no bias}
% field-shifted edge (red-shifted to near 1545 nm, exciton broadened)
\addplot[engred] coordinates {
  (1490,1.00) (1500,0.95) (1515,0.85) (1525,0.72) (1535,0.55)
  (1542,0.40) (1548,0.28) (1552,0.20) (1558,0.10) (1565,0.045)
  (1572,0.02) (1580,0.01)
};
\addlegendentry{reverse bias}
% operating wavelength
\addplot[enggray, dashed, thick] coordinates {(1550,0) (1550,1.2)};
\node[enggray, anchor=south, font=\tiny, rotate=90]
  at (axis cs:1550,0.55) {operating $\lambda$};
% mark the two states at operating wavelength
\addplot[engblue, only marks, mark=*, mark size=1.6pt]
  coordinates {(1550,0.005)};
\addplot[engred, only marks, mark=square*, mark size=1.8pt]
  coordinates {(1550,0.15)};
\node[engblue, anchor=west, font=\tiny] at (axis cs:1552,0.02) {on};
\node[engred, anchor=west, font=\tiny] at (axis cs:1552,0.17) {off};
\end{axis}
\end{tikzpicture}
\caption[Electro-absorption modulator edge shift]{Absorption edge of a
multiple-quantum-well modulator at zero bias and under reverse bias. The applied
field pulls the electron and hole wavefunctions to opposite sides of each well and
lowers the transition energy, shifting the excitonic absorption edge to longer
wavelength, the quantum-confined Stark effect. At a fixed operating wavelength
chosen just beyond the zero-field edge the device is nearly transparent with no
bias and strongly absorbing when the bias is applied, so a small voltage switches
the transmitted light between an on and an off state. The extinction ratio is the
ratio of the two transmissions, and the modulator's speed is set by the tiny
capacitance of the reverse-biased well stack rather than by any carrier
lifetime.}
\label{fig:vol04ch12:eam-absorption}
\end{figure}
